\section{Problem formulation and privacy model}
\label{sec:privacy}
\subsection{Personalized implicit-feedback ranking}
Let $\mathcal I_u^+$ be user $u$'s set of training-positive items. A local vector
$p_u\in\R^d$ and shared item row $q_i\in\R^d$ produce score
\begin{equation}
 s_{ui}=p_u^Tq_i,\qquad S=PQ^T.
 \label{eq:score}
\end{equation}
The model has no item biases, content features, attention layers, or temporal
features. For a sampled triplet $(u,i,j)$ with $i\in\mathcal I_u^+$ and
$j\notin\mathcal I_u^+$, the implemented BPR loss is
\begin{equation}
 \ell(u,i,j)=\log(1+\exp(s_{uj}-s_{ui}))+
 \lambda\bigl(\norm{p_u}^2+\norm{q_i}^2+\norm{q_j}^2\bigr).
 \label{eq:bpr}
\end{equation}
The softplus implementation is numerically stable for extreme margins. The
penalty applies to embeddings touched by the sampled loss, not to a separate
factor-only penalty or an optimizer's implicit weight decay. In a factorized
model, $q_i=A_iB$ and gradients propagate through this effective row.

There are two relevant population objectives. Interaction-weighted centralized
training gives more weight to users with more positives. User-weighted training
first averages within each user and then across users. Uniform client inclusion
and equally weighted client updates are closer to the latter, although local
multi-step optimization is not exactly the gradient of a single centralized
objective. The B0-UW control makes the weighting difference visible without
claiming that it eliminates all optimizer differences.

\subsection{Protected information and adjacency}
The privacy unit is one user's complete interaction record under add/remove
adjacency. For neighboring datasets $D,D'$ differing in one user's contribution,
a mechanism $\mathcal M$ is $(\epsilon,\delta)$-DP when, for each measurable
output set $\mathcal O$,
\begin{equation}
 \Pr[\mathcal M(D)\in\mathcal O]
 \leq e^\epsilon\Pr[\mathcal M(D')\in\mathcal O]+\delta.
 \label{eq:dp}
\end{equation}
The catalog, eligible population size used for normalization, configurations,
and random public basis are treated as fixed public inputs. Eligibility filtering
and catalog construction in the research data pipeline are not themselves
privatized. Thus the reported mechanism is conditional on these inputs; it is
not an end-to-end private data-ingestion and model-selection system.

User vectors persist locally and are never included in the shared update.
The trust model assumes a server that sees only a noised aggregate, as would
be supported by suitable secure aggregation and noise generation. The simulator
does not implement a cryptographic secure-aggregation protocol. Raw client
updates or local checkpoints visible to an untrusted server would not be
protected merely by calling the training procedure federated.

\subsection{Private aggregation mechanism}
Write $\theta_t$ for the shared parameter vector: vectorized $Q$, $(A,B)$, or
$A$ alone. At round $t$, each client is independently sampled with probability
$q$. A selected client starts from $\theta_t$, performs local updates using its
own data and local $p_u$, and returns shared difference $\Delta_{u,t}$. Clip
the whole shared contribution once:
\begin{equation}
 \bar\Delta_{u,t}=\Delta_{u,t}
 \min\left(1,\frac{C}{\norm{\Delta_{u,t}}_2}\right),\qquad
 \theta_{t+1}=\theta_t+
 \frac{\eta_s}{qN}\left(\sum_{u\in\mathcal U_t}\bar\Delta_{u,t}
            +\xi_t\right),
 \quad\xi_t\sim\mathcal N(0,\sigma^2C^2I).
 \label{eq:aggregate}
\end{equation}
For Two, the norm is over the concatenated factor update. Clipping each factor
independently at $C$ would define a different sensitivity bound. Noise is added
to every shared coordinate, including item rows untouched by sampled users.
The denominator is the fixed expectation $qN$, not the realized number of
participants. Empty private rounds still release a noised update and count
toward composition. Local states are retained only by their respective users.

At the ML-100K population size, $qN=94.2$. At ML-1M, $qN=603.5$. The primary
accountant depends on $q,T,\sigma,\delta$, so transferring these parameters
preserves the accountant bound while changing the post-normalization noise SD.
Changing dimension alone does not reduce $\epsilon$. Neither does changing
$C$ while scaling Gaussian noise by the same multiplier. Utility can change
substantially in both cases.

% Generated by scripts/generate_tables.py; do not edit numeric cells.
\begin{table}[tbp]\centering\small
\caption{Final iterative privacy calibration at $q=0.1$, $T=50$, $\delta=10^{-5}$.}\label{tab:privacy}
\setlength{\tabcolsep}{4pt}
\begin{tabular}{lrrr}\toprule
$\epsilon$ target & $\sigma$ & PRV bound & RDP cross-check \\
\midrule
8 & 0.805072021 & 7.967513 & 9.119818 \\
4 & 1.151928711 & 3.963885 & 4.497283 \\
2 & 1.760449219 & 1.993278 & 2.220645 \\
1 & 2.977490234 & 0.990062 & 1.087196 \\
\bottomrule\end{tabular}
\tabnote{PRV is primary. RDP is a separate bound on the same mechanism. Source: privacy\_accounting.csv, final 50-round rows.}
\end{table}

All final private collaborative experiments use $T=50$, $q=0.1$, $C=1$, and
$\delta=10^{-5}$. \Cref{tab:privacy} reports the exact retained calibration.
PRV is the declared primary bound. The RDP column is a distinct upper bound
for the same mechanism, not an extra privacy charge and not a replacement
selected after inspecting utility.

\subsection{Release scope and local personalization}
The guarantee describes one fixed-hyperparameter training mechanism. Publishing
multiple runs requires composition; selecting a configuration from private
validation results is not free; reporting user-level diagnostics is not covered
by the final aggregate's guarantee. The research process uses nonprivate
validation, evaluation, and access to internal states. Reproducible random
seeds are useful for auditing but do not constitute production-grade secure
randomness.

Local personalization can be interpreted through the established joint-DP
framework when a user's output depends on that user's own records and an
appropriately private shared transcript. This does not authorize publishing
the entire local matrix $P$. E2 and E3 use oracle reset and alignment operations
on private states for diagnosis; E4 conditions on private checkpoints. They
are research measurements, not releases justified by the starting checkpoint's
privacy label. These distinctions prevent the formal accountant from being
mistaken for an end-to-end privacy claim about the report itself.
